\documentclass[a5paper,10pt]{book}

\usepackage{amsfonts}
\usepackage{hyperref}
\usepackage[utf8]{inputenc}
\usepackage[dvipsnames]{xcolor}
\usepackage[hmarginratio=1:1]{geometry}% equal left and right margins
\usepackage{graphicx}
% \usepackage{fontspec}
% \setmainfont{Times New Roman}
\usepackage{titling}

% \usepackage{sectsty}
% \chapterfont{\color{Red} \large}  % sets colour of chapters
% \sectionfont{\color{Red}}  % sets colour of sections

\usepackage{afterpage}
\newcommand\myemptypage{
    \null
    \thispagestyle{empty}
    \addtocounter{page}{-1}
    \newpage
    }

\setlength{\droptitle}{-1in}

\def\ae{æ}
\def\AE{Æ}



\begin{document}
\pagenumbering{roman}


\title{\includegraphics[width=2in]{Picture1.png} \\ \LARGE \bfseries \color{Red} T H E \\ \Huge R O M A N \\ \huge B R E V I A R Y, \\ \large \mdseries \color{Black} From the most sacred scripture, \& proven by the \\ newly acquired histories of the saints, and again \\ more accurately revised by the \\ same author. \\ \ \\ \itshape Search the scriptures, \\ for they give testimony of me. \\John 5:39}

\author{PRINTED IN VENICE}
\date{MDCXXXI}

\maketitle

\myemptypage

% \frontmatter

\pagenumbering{roman}

\begin{center} \color{red}
AD SANCTISSIMVM PATREM \\ \& Dominum nostrum Paulum tertium \\ Pontificem Maximum, \\ \color{black} Francisci Quignonij Tituli Sanctæ Crucis in Ierusalem presbyteri Cardinalis, in Breuiarium proxime consectum, ac denuo \\ recognitum.
\end{center}

\centerline{\color{red}P R \AE \ F A T I O.}

\Huge \textbf{W}
% \normalsize Reviarum Romanum nuper a nobis felicis recordationis Clementis, vii. Pontifi. Maximi hortatu consectum, ac potius in ampliorem sacrarum scripturam lectionem, ad veterem sanctorum patrum, \& conciliorum antiquorum formam reuocatum, tuaque voluntate sanctissime pater editum, graues plerosque ac doctos viros ita probasse \& recepisse intellexi, vt nihil in eo metandum existimarent.
\normalsize E present the Roman Breviary by the urging of the Supreme Pontiff Clement VII of happy memory, and more importantly, for a wider reading of sacred scripture, to the holy fathers, \& to restore the form of the ancient councils, published by your will most holy father, and as I understand, the most serious to educated men approved and received it, that they thought there was no need to stop there.
% Alios item animaduerti graues etiam, \& prudentes homines, qui eius rationem magnopere probantes, nonnihil tamen in eo desiderari affimarent.
Other serious items to be noticed also, \& prudentes homines, qui eius rationem magnopere probantes, nonnihil tamen in eo desiderari affimarent.
Illud vero nunquam dubitatui, fore in tanta multitudine nonnullos ex ijs videlicet, qui in diuerso precandi ritu consenuissent, quibus labor ille noster non esset perinde gratus, existimantibus ab inueterata illa consuetudine precandi nulla ratione clericis esse discedendum.
Imo vero nobis primam editionem Breuiarij, non tamquam promulgationem legis esse placuerat, sed quasi publicam quandam deliberationem:
vt sic proposita nostra sententia, iudicia multorum exquireremus, \& quod omnium commodissimum, \& religioni, ac pietati conuenientissimum plerisque prudentibus, grauibusque viris visum esset, sequeremur.
Nihil enim humano elaboratum ingenio tam exactum initio vnquam fuit, quin postea multorum accedente iudicio perfectius reddi possit, vt in ipsis etiam ecclesiasticis institutis, circa primitiuam presertim ecclesiam contigisse videmus.
Itaque, multorum sententiis % end pp8
collatis, qu\ae \ nobis partim vocibus, partim scriptis innotuerunt, iudicium eorum secuti, qui omnium prudentissime sentire visi sunt, libenter quedam addidimus, alia mutauimus, \& omnia dilidenter recongnouimus, retenta tamen summa forma Breuiarii.
Sed quoniam sic fert natura rerum, vt nihil sit tam rectum, nihil ranta ratione in vitam vsumque hominum inductum, cuius nouitas non sit aliquibus ingrata, non temere facturi esse videmur, si rationem totius instituti nostri a nobis prius summatim redditam, nunc accuratius recognito Breuiario, paulo latius explicabimus.
Mihi enim (vt s\ae pe sum professus) cogitanti, atque animo repetenti initia veteris insituti, quo sancitum est, vt clerici sacris initiati, vel sacerdotiis pr\ae sidentes, singulis diebus perlegant horarias preces, quas Canonicas etiam appellamus, tres pr\ae cipue caus\ae \ spectat\ae \ fuisse videri solent.
Quarum ea prima est, qui cum c\ae teri homines in quaque ciuitate, aut suum quisque negocium agant, aut in Republica administranda sint occupati, clericis ex eo vocatis, qui sortis Domini sint, quique bonis ecclesiasticis aluntur, hoc potissimum negocium diuinis, \& humanis legibus est iniunctum, vt Deum habere propitium in comissum sibi populum, de se bene merentem cunctis rationibus enitantur.
quod %sic
non solum sacrificiis efficitur, sed etiam precibus, qu\ae \ a pio corede proficiscantur.
ipsi %sic
enim (vt testimonio versibusque vtar beati Alexandri pap\ae \ \& martyris) pro populo interpellant, \& populi peccata comedunt, quia precibus suis \& oblationibus ea delent, atque consumunt.
Qui quanto

\chapter{Fundamentals of Vectors}

\section{Vectors in Euclidean Space}

% \mainmatter

\pagenumbering{arabic}

\chapter{First Chapter}

This will be an empty chapter...

\newpage
\thispagestyle{empty}
\mbox{}
\newpage

\chapter{Fundamentals of Vectors}

\section{Vectors in Euclidean Space}







\end{document}

